\documentclass[12pt,a4paper]{article}
\usepackage[utf8]{inputenc}
\usepackage[T1]{fontenc}
\usepackage[french]{babel}
\usepackage{amsmath,amssymb}
\usepackage{geometry}
\geometry{margin=1.8cm}
\usepackage{enumitem}
\usepackage{parskip}
\usepackage{array}
\usepackage{booktabs}
\usepackage{eurosym}
\usepackage{xcolor}
\usepackage{tcolorbox}
\tcbuselibrary{skins,breakable}
\usepackage{tikz}
\usepackage{multicol}
\usepackage{pifont}

% --- Couleurs du jeu ---
\definecolor{bleufonce}{RGB}{20,60,120}
\definecolor{rouge}{RGB}{180,30,30}
\definecolor{vert}{RGB}{20,130,70}
\definecolor{orange}{RGB}{230,120,20}
\definecolor{violet}{RGB}{110,30,140}
\definecolor{grisclair}{RGB}{240,240,240}
\definecolor{or}{RGB}{220,170,20}

% --- Boîtes stylisées ---
\newtcolorbox{regle}[1][]{
  colback=grisclair, colframe=bleufonce, boxrule=1pt,
  arc=3mm, left=3mm, right=3mm, top=2mm, bottom=2mm, #1
}

\newtcolorbox{carte}[2]{
  colback=white, colframe=#1, boxrule=1.2pt,
  arc=2mm, left=3mm, right=3mm, top=2mm, bottom=2mm,
  title=#2, fonttitle=\bfseries
}

\newtcolorbox{corrige}[1][]{
  colback=vert!5, colframe=vert, boxrule=1pt,
  arc=2mm, left=3mm, right=3mm, top=2mm, bottom=2mm,
  title=Corrigé, fonttitle=\bfseries, #1
}

\newtcolorbox{conseil}[1][]{
  colback=orange!5, colframe=orange, boxrule=1pt,
  arc=2mm, left=3mm, right=3mm, top=2mm, bottom=2mm,
  title=Conseil d'animation, fonttitle=\bfseries, #1
}

% --- Étoiles pour le barème progressif ---
\newcommand{\etoiles}[1]{%
  \ifcase#1
  \or \textcolor{or}{\ding{72}}%
  \or \textcolor{or}{\ding{72}\ding{72}}%
  \or \textcolor{or}{\ding{72}\ding{72}\ding{72}}%
  \or \textcolor{or}{\ding{72}\ding{72}\ding{72}\ding{72}}%
  \fi
}

\begin{document}

\begin{center}
{\Huge\bfseries\color{bleufonce} Le Grand Défi des Évolutions}\\[0.3cm]
{\Large\bfseries\color{rouge} Version enseignant -- Corrigés et conseils}\\[0.3cm]
{\large Jeu mathématique -- Durée : 55 minutes}\\[0.2cm]
{\large Niveau : Première / Terminale -- Taux, coefficients multiplicateurs, taux global et moyen}
\end{center}

\vspace{0.3cm}

% ============================================================
\section*{Organisation générale}

\begin{regle}
\textbf{Objectifs pédagogiques :}
\begin{itemize}
    \item Consolider la notion de coefficient multiplicateur.
    \item Distinguer taux global et taux moyen.
    \item Comprendre la non-additivité des taux.
    \item Développer la persévérance et le travail en équipe.
\end{itemize}
\textbf{Modalités :}
\begin{itemize}
    \item 4 équipes de 6 à 8 élèves, un capitaine par équipe.
    \item 5 manches chronométrées, un tableau de score affiché ou tenu par l'enseignant.
    \item Calculatrice autorisée.
\end{itemize}
\end{regle}

\vspace{0.3cm}

\begin{center}
\renewcommand{\arraystretch}{1.4}
\begin{tabular}{|c|l|c|c|}
\hline
\textbf{Manche} & \textbf{Titre} & \textbf{Durée} & \textbf{Points max} \\
\hline
1 & Échauffement éclair & 10 min & 8 pts \\
\hline
2 & Défis en équipe & 15 min & 12 pts \\
\hline
3 & Duel de capitaines & 10 min & 8 pts \\
\hline
4 & Enveloppe mystère & 10 min & 12 pts \\
\hline
5 & Finale « Tout ou rien » & 10 min & 10 pts \\
\hline
\textbf{Total} & & \textbf{55 min} & \textbf{50 pts} \\
\hline
\end{tabular}
\end{center}

\newpage

% ============================================================
\section*{Manche 1 -- Échauffement éclair (10 min, 8 pts)}

\textit{Chaque bonne réponse rapporte 1 point. Les 8 questions sont posées à toute la classe. La première équipe qui lève la main et donne la bonne réponse marque le point.}

\vspace{0.3cm}

\begin{carte}{bleufonce}{Question 1}
Un prix augmente de \(25\%\). Quel est le coefficient multiplicateur ?
\end{carte}
\begin{corrige}
\(CM = 1 + \dfrac{25}{100} = 1{,}25\).
\end{corrige}

\begin{carte}{bleufonce}{Question 2}
Un prix baisse de \(10\%\). Quel est le coefficient multiplicateur ?
\end{carte}
\begin{corrige}
\(CM = 1 - \dfrac{10}{100} = 0{,}90\).
\end{corrige}

\begin{carte}{bleufonce}{Question 3}
Le coefficient multiplicateur est \(1{,}4\). Quel est le taux d'évolution en \% ?
\end{carte}
\begin{corrige}
\(t = (1{,}4 - 1) \times 100 = +40\%\).
\end{corrige}

\begin{carte}{bleufonce}{Question 4}
Le coefficient multiplicateur est \(0{,}75\). Quel est le taux d'évolution en \% ?
\end{carte}
\begin{corrige}
\(t = (0{,}75 - 1) \times 100 = -25\%\).
\end{corrige}

\begin{carte}{bleufonce}{Question 5}
Un article passe de \(50\)~\euro{} à \(60\)~\euro{}. Quel est le taux d'évolution ?
\end{carte}
\begin{corrige}
\(t = \dfrac{60-50}{50} \times 100 = \dfrac{10}{50} \times 100 = +20\%\).
\end{corrige}

\begin{carte}{bleufonce}{Question 6}
Une grandeur augmente de \(20\%\), puis de \(20\%\). Le taux global est-il \(40\%\) ?
\end{carte}
\begin{corrige}
Non. \(CM_{\text{global}} = 1{,}20 \times 1{,}20 = 1{,}44\), soit \(t = +44\%\). Les taux ne s'additionnent pas : ils se composent.
\end{corrige}

\begin{carte}{bleufonce}{Question 7}
Que vaut \(1{,}10 \times 0{,}90\) ?
\end{carte}
\begin{corrige}
\(1{,}10 \times 0{,}90 = 0{,}99\). Une hausse de \(10\%\) suivie d'une baisse de \(10\%\) donne une baisse globale de \(1\%\).
\end{corrige}

\begin{carte}{bleufonce}{Question 8}
Un prix baisse de \(50\%\). Par quel coefficient faut-il multiplier pour revenir au prix initial ?
\end{carte}
\begin{corrige}
Il faut multiplier par \(\dfrac{1}{0{,}5} = 2\). Vérification : \(0{,}5 \times 2 = 1\).
\end{corrige}

\begin{conseil}
Annoncez les questions rapidement (30 secondes max par question). Utilisez un buzzer ou un simple lever de main. Faites reformuler la réponse par un autre membre de l'équipe pour vérifier la compréhension.
\end{conseil}

\newpage

% ============================================================
\section*{Manche 2 -- Défis en équipe (15 min, 12 pts)}

\textit{Chaque équipe reçoit les 3 défis ci-dessous. Elle dispose de 15 minutes pour les résoudre. Chaque défi est noté sur 4 points (2 pts pour la démarche, 2 pts pour le résultat).}

\vspace{0.3cm}

\begin{carte}{rouge}{Défi A}
Un smartphone coûte \(600\)~\euro{}. Son prix baisse de \(15\%\), puis augmente de \(10\%\).\\[0.2cm]
\textbf{Quel est son prix final ? Le prix a-t-il globalement augmenté ou baissé ?}
\end{carte}
\begin{corrige}
\(CM_1 = 1 - 0{,}15 = 0{,}85\).\quad \(CM_2 = 1 + 0{,}10 = 1{,}10\).\\[0.2cm]
\(CM_{\text{global}} = 0{,}85 \times 1{,}10 = 0{,}935\).\\[0.2cm]
\(V_f = 600 \times 0{,}935 = 561\)~\euro{}.\\[0.2cm]
Le coefficient global étant inférieur à \(1\), le prix a globalement \textbf{baissé de \(6{,}5\%\)}.
\end{corrige}

\begin{carte}{rouge}{Défi B}
Une ville comptait \(25\,000\) habitants en 2015. Elle a perdu \(4\%\) d'habitants en 2016, puis gagné \(6\%\) en 2017, puis perdu \(2\%\) en 2018.\\[0.2cm]
\textbf{Combien d'habitants compte-t-elle fin 2018 (arrondir à l'unité) ? Quel est le taux global d'évolution ?}
\end{carte}
\begin{corrige}
\(CM_{\text{global}} = 0{,}96 \times 1{,}06 \times 0{,}98 = 0{,}997248\).\\[0.2cm]
\(V_f = 25\,000 \times 0{,}997248 = 24\,931{,}2 \approx 24\,931\) habitants.\\[0.2cm]
\(t = (0{,}997248 - 1) \times 100 = -0{,}2752\%\).\\[0.2cm]
La population a très légèrement baissé.
\end{corrige}

\begin{carte}{rouge}{Défi C}
Un article subit deux évolutions successives : \(+30\%\) puis \(-30\%\).\\[0.2cm]
\textbf{Le prix final est-il égal au prix initial ? Justifier par un calcul.}
\end{carte}
\begin{corrige}
\(CM_{\text{global}} = 1{,}30 \times 0{,}70 = 0{,}91 \neq 1\).\\[0.2cm]
Le prix final n'est pas égal au prix initial : il est \textbf{inférieur de \(9\%\)}.\\[0.2cm]
Piège classique : \(+30\%\) puis \(-30\%\) ne s'annule pas car les pourcentages ne portent pas sur la même valeur de référence.
\end{corrige}

\begin{conseil}
Circulez entre les équipes. Vérifiez que les élèves passent bien par les coefficients multiplicateurs et non par des additions de taux. Exigez une phrase de conclusion.
\end{conseil}

\newpage

% ============================================================
\section*{Manche 3 -- Duel de capitaines (10 min, 8 pts)}

\textit{Les 4 capitaines s'affrontent. Chaque bonne réponse rapporte 2 points à son équipe. Les autres membres peuvent aider mais seul le capitaine répond.}

\vspace{0.3cm}

\begin{carte}{orange}{Duel 1}
Un capital augmente de \(10\%\) la première année, puis de \(20\%\) la deuxième.\\[0.2cm]
\textbf{Quel taux unique aurait donné le même résultat en 2 ans ?}
\end{carte}
\begin{corrige}
\(CM_{\text{global}} = 1{,}10 \times 1{,}20 = 1{,}32\).\\[0.2cm]
\(t = (1{,}32 - 1) \times 100 = +32\%\).\\[0.2cm]
Un taux unique de \(+32\%\) sur 2 ans donne le même résultat.
\end{corrige}

\begin{carte}{orange}{Duel 2}
Un prix a augmenté de \(44\%\) en deux ans, avec le même taux annuel \(t\).\\[0.2cm]
\textbf{Déterminer \(t\).}
\end{carte}
\begin{corrige}
\((1+t)^2 = 1{,}44\), donc \(1+t = \sqrt{1{,}44} = 1{,}2\).\\[0.2cm]
\(t = 0{,}2\), soit \(\boxed{t = 20\%}\).
\end{corrige}

\begin{carte}{orange}{Duel 3}
Une grandeur a subi 3 évolutions successives : \(+10\%\), \(-20\%\), \(+25\%\).\\[0.2cm]
\textbf{Quel est le taux global ?}
\end{carte}
\begin{corrige}
\(CM_{\text{global}} = 1{,}10 \times 0{,}80 \times 1{,}25 = 1{,}10\).\\[0.2cm]
\(t = (1{,}10 - 1) \times 100 = +10\%\).
\end{corrige}

\begin{carte}{orange}{Duel 4}
Un prix passe de \(80\)~\euro{} à \(100\)~\euro{}.\\[0.2cm]
\textbf{Quel est le taux d'évolution ? Puis, quel taux faudrait-il appliquer pour revenir à \(80\)~\euro{} ?}
\end{carte}
\begin{corrige}
Hausse : \(t_1 = \dfrac{100-80}{80} \times 100 = +25\%\).\\[0.2cm]
Retour : \(t_2 = \dfrac{80-100}{100} \times 100 = -20\%\).\\[0.2cm]
On vérifie : \(1{,}25 \times 0{,}80 = 1\).
\end{corrige}

\begin{conseil}
Faites répondre uniquement le capitaine. Les autres membres peuvent souffler mais pas écrire. Chronométrez 2 minutes max par duel pour garder le rythme.
\end{conseil}

\newpage

% ============================================================
\section*{Manche 4 -- Enveloppe mystère (10 min, 12 pts)}

\begin{regle}
\textbf{Le défi de la persévérance !}\\[0.2cm]
Chaque équipe démarre avec une enveloppe. Dès qu'elle a terminé, elle la rend et reçoit l'enveloppe suivante (rotation dans le sens horaire).\\[0.2cm]
\textbf{Objectif :} en résoudre le plus possible en 10 minutes.\\[0.3cm]
\textbf{Barème progressif :}\\[0.3cm]
\begin{center}
\renewcommand{\arraystretch}{1.6}
\begin{tabular}{|>{\raggedright\arraybackslash}p{5.5cm}|c|c|}
\hline
\textbf{Enveloppe résolue} & \textbf{Étoiles} & \textbf{Points} \\
\hline
1\textsuperscript{re} enveloppe correcte & \etoiles{1} & \(2\) pts \\
\hline
2\textsuperscript{e} enveloppe correcte & \etoiles{2} & \(3\) pts \\
\hline
3\textsuperscript{e} enveloppe correcte & \etoiles{3} & \(3\) pts \\
\hline
4\textsuperscript{e} enveloppe correcte & \etoiles{4} & \(4\) pts \\
\hline
\multicolumn{2}{|r|}{\textbf{Total maximal}} & \(\mathbf{12}\) \textbf{pts} \\
\hline
\end{tabular}
\end{center}
\vspace{0.2cm}
\textbf{Règle d'or :} une enveloppe non résolue ou fausse ne rapporte rien (pas de malus).\\[0.2cm]
\textbf{Astuce :} ne restez pas bloqués ! Si une enveloppe résiste, rendez-la et passez à la suivante.
\end{regle}

\vspace{0.5cm}

\begin{conseil}
Prévoyez plusieurs copies de chaque enveloppe (au moins 2 par équipe) pour fluidifier la rotation. Numérotez les enveloppes au dos pour suivre l'ordre de passage.
\end{conseil}

\newpage

% --- Corrigés des 4 enveloppes ---

\subsection*{Enveloppe A -- Corrigé}

\begin{carte}{violet}{Enveloppe A}
Un commerçant augmente ses prix de \(20\%\). Pour compenser une baisse des ventes, il les diminue ensuite de \(20\%\).\\[0.3cm]
\textbf{Les prix sont-ils revenus à leur niveau initial ? Justifier.}
\end{carte}
\begin{corrige}
\(CM_{\text{global}} = 1{,}20 \times 0{,}80 = 0{,}96 \neq 1\).\\[0.2cm]
Non, les prix ne sont pas revenus à leur niveau initial : ils ont \textbf{baissé de \(4\%\)} par rapport au prix de départ.\\[0.2cm]
Raison : la baisse de \(20\%\) s'applique sur un prix déjà augmenté de \(20\%\), donc sur une valeur de référence plus grande.
\end{corrige}

\subsection*{Enveloppe B -- Corrigé}

\begin{carte}{violet}{Enveloppe B}
Un article coûte \(250\)~\euro{} après une augmentation de \(25\%\).\\[0.3cm]
\textbf{Quel était son prix avant l'augmentation ?}
\end{carte}
\begin{corrige}
\(V_f = CM \times V_i \Rightarrow V_i = \dfrac{V_f}{CM}\).\\[0.2cm]
\(CM = 1 + \dfrac{25}{100} = 1{,}25\).\\[0.2cm]
\(V_i = \dfrac{250}{1{,}25} = 200\)~\euro{}.\\[0.2cm]
Vérification : \(200 \times 1{,}25 = 250\).
\end{corrige}

\subsection*{Enveloppe C -- Corrigé}

\begin{carte}{violet}{Enveloppe C}
Un placement a rapporté \(21\%\) en deux ans, avec le même taux annuel \(t\).\\[0.3cm]
\textbf{Déterminer \(t\) (arrondir à \(0{,}01\%\)).}
\end{carte}
\begin{corrige}
\(CM_{\text{global}} = 1 + \dfrac{21}{100} = 1{,}21\).\\[0.2cm]
\((1+t)^2 = 1{,}21 \Rightarrow 1+t = \sqrt{1{,}21} = 1{,}10\).\\[0.2cm]
\(t = 10\%\). Le taux annuel moyen est de \(10\%\).
\end{corrige}

\subsection*{Enveloppe D -- Corrigé}

\begin{carte}{violet}{Enveloppe D}
Une population augmente de \(5\%\) par an pendant 3 ans.\\[0.3cm]
\textbf{Quel est le taux global ? Quel est le taux annuel moyen ?}
\end{carte}
\begin{corrige}
\(CM_{\text{global}} = 1{,}05^3 = 1{,}157625\).\\[0.2cm]
\(t_{\text{global}} = (1{,}157625 - 1) \times 100 = +15{,}76\%\) (arrondi à \(0{,}01\%\)).\\[0.2cm]
Le taux annuel moyen est \(t_m = 5\%\) (car la croissance est déjà à taux constant).\\[0.2cm]
Vérification : \((1{,}05)^3 = 1{,}157625\).
\end{corrige}

\newpage

% ============================================================
\section*{Manche 5 -- Finale « Tout ou rien » (10 min, 10 pts)}

\textit{Chaque équipe choisit un niveau de difficulté. Une bonne réponse rapporte les points indiqués ; une mauvaise réponse les retire. Les points peuvent devenir négatifs.}

\vspace{0.3cm}

\begin{center}
\renewcommand{\arraystretch}{1.5}
\begin{tabular}{|c|p{8.5cm}|c|c|}
\hline
\textbf{Niveau} & \textbf{Question} & \textbf{Gain} & \textbf{Perte} \\
\hline
Facile & Un prix augmente de \(30\%\). Quel coefficient appliquer pour revenir au prix initial ? (arrondir à \(0{,}001\)) & \(+2\) & \(-2\) \\
\hline
Moyen & Un capital de \(5\,000\)~\euro{} subit 4 évolutions : \(+8\%\), \(-5\%\), \(+12\%\), \(-10\%\). Quel est le capital final ? & \(+3\) & \(-3\) \\
\hline
Difficile & Une grandeur a subi 5 évolutions annuelles identiques de taux \(t\). Le taux global est \(+27{,}628\%\). Déterminer \(t\). & \(+5\) & \(-5\) \\
\hline
\end{tabular}
\end{center}

\vspace{0.3cm}

\begin{corrige}[title={Corrigé -- Niveau facile}]
Hausse de \(30\%\) : \(CM = 1{,}30\).\\[0.2cm]
Coefficient de retour : \(\dfrac{1}{1{,}30} \approx 0{,}769\).\\[0.2cm]
Soit un taux de retour d'environ \(-23{,}1\%\).
\end{corrige}

\begin{corrige}[title={Corrigé -- Niveau moyen}]
\(CM_{\text{global}} = 1{,}08 \times 0{,}95 \times 1{,}12 \times 0{,}90\).\\[0.2cm]
\(1{,}08 \times 0{,}95 = 1{,}026\).\quad \(1{,}026 \times 1{,}12 = 1{,}14912\).\quad \(1{,}14912 \times 0{,}90 = 1{,}034208\).\\[0.2cm]
\(V_f = 5\,000 \times 1{,}034208 = 5\,171{,}04\)~\euro{}.
\end{corrige}

\begin{corrige}[title={Corrigé -- Niveau difficile}]
\((1+t)^5 = 1{,}27628\).\\[0.2cm]
\(1+t = (1{,}27628)^{1/5} = 1{,}05\).\\[0.2cm]
\(t = 5\%\).
\end{corrige}

\begin{conseil}
Annoncez les points avant la réponse pour créer du suspense. Rappelez que les points peuvent devenir négatifs : cela incite les équipes à bien choisir leur niveau. En cas d'égalité finale, posez la question bonus : \textit{« Un prix augmente de \(100\%\), puis baisse de \(50\%\). Revient-il au prix initial ? »} (Réponse : oui, car \(2 \times 0{,}5 = 1\).)
\end{conseil}

\newpage

% ============================================================
\section*{Fiche de score (version enseignant)}

\begin{center}
{\large \textbf{Nom de l'équipe :} \dotfill}\\[0.3cm]
{\large \textbf{Capitaine :} \dotfill}
\end{center}

\vspace{0.5cm}

\begin{center}
\renewcommand{\arraystretch}{1.8}
\begin{tabular}{|c|c|c|}
\hline
\textbf{Manche} & \textbf{Points max} & \textbf{Points obtenus} \\
\hline
1 -- Échauffement éclair & 8 & \\
\hline
2 -- Défis en équipe & 12 & \\
\hline
3 -- Duel de capitaines & 8 & \\
\hline
4 -- Enveloppe mystère & 12 & \\
\hline
5 -- Finale « Tout ou rien » & 10 & \\
\hline
\textbf{Total} & \textbf{50} & \\
\hline
\end{tabular}
\end{center}

\vspace{0.6cm}

\section*{Détail Manche 4 -- Enveloppe mystère}

\begin{center}
\renewcommand{\arraystretch}{1.6}
\begin{tabular}{|c|c|c|c|c|}
\hline
\textbf{Ordre} & \textbf{Enveloppe} & \textbf{Correcte ? (O/N)} & \textbf{Étoiles} & \textbf{Points} \\
\hline
1 & & & \etoiles{1} & /2 \\
\hline
2 & & & \etoiles{2} & /3 \\
\hline
3 & & & \etoiles{3} & /3 \\
\hline
4 & & & \etoiles{4} & /4 \\
\hline
\multicolumn{4}{|r|}{\textbf{Total Manche 4}} & \textbf{/12} \\
\hline
\end{tabular}
\end{center}

\vspace{0.6cm}

\section*{Récapitulatif des notions mobilisées}

\begin{center}
\renewcommand{\arraystretch}{1.5}
\begin{tabular}{|>{\raggedright\arraybackslash}p{7cm}|c|}
\hline
\textbf{Notion} & \textbf{Manches concernées} \\
\hline
Coefficient multiplicateur \(CM = 1 + \dfrac{t}{100}\) & 1, 2, 3, 4, 5 \\
\hline
Relation \(V_f = CM \times V_i\) & 1, 2, 3, 4 \\
\hline
Évolutions successives : \(CM_{\text{global}} = CM_1 \times CM_2 \times \cdots\) & 2, 3, 4, 5 \\
\hline
Taux global \(T = (CM_{\text{global}} - 1) \times 100\) & 2, 3, 4, 5 \\
\hline
Taux moyen : \((1+t_m)^n = CM_{\text{global}}\) & 3, 4, 5 \\
\hline
Non-additivité des taux & 1, 2, 3, 4 \\
\hline
Résolution d'équation avec puissance/racine & 3, 4, 5 \\
\hline
\end{tabular}
\end{center}

\vspace{0.6cm}

\begin{conseil}
\textbf{Bilan de fin de séance :} prenez 3 à 5 minutes pour faire verbaliser les élèves sur les pièges rencontrés (additivité des taux, valeur de référence, oubli du coefficient de retour). Ce temps de métacognition est essentiel pour ancrer les notions.
\end{conseil}

\vspace{0.8cm}

\begin{center}
{\Large\bfseries\color{bleufonce} Bon jeu, et bonne évaluation formative !}
\end{center}
\newpage

% ============================================================
\section*{Mots-clés et symboles pour l'observation rapide}

\begin{regle}
	\textbf{Principe :} pendant le jeu, l'enseignant ne rédige pas. Il note \textbf{un symbole} ou \textbf{un mot-clé} par équipe et par manche (5 secondes max). Ces notes servent ensuite à remplir la fiche « Observations par équipe » après la séance, en moins de 5 minutes.\\[0.2cm]
	\textbf{Objectif :} garder une trace fine et exploitable, sans ralentir l'animation.
\end{regle}

\vspace{0.3cm}

\subsection*{Code abrégé -- Points forts}

\begin{center}
	\renewcommand{\arraystretch}{1.4}
	\begin{tabular}{|c|>{\raggedright\arraybackslash}p{6cm}|>{\raggedright\arraybackslash}p{7.5cm}|}
		\hline
		\textbf{Symbole} & \textbf{Signification} & \textbf{Exemple observé} \\
		\hline
		\texttt{CM} & Maîtrise le passage taux \(\to\) coefficient multiplicateur & « \(1{,}15\) trouvé immédiatement » \\
		\hline
		\texttt{MULT} & Multiplie correctement les coefficients (évolutions successives) & « \(0{,}94 \times 1{,}10 \times 0{,}96\) posé sans hésiter » \\
		\hline
		\texttt{VF} & Applique correctement \(V_f = CM \times V_i\) & « \(600 \times 0{,}935 = 561\) » \\
		\hline
		\texttt{INT} & Interprète correctement le signe du résultat & « baisse de \(6{,}5\%\) » dit spontanément \\
		\hline
		\texttt{RAC} & Utilise la racine n-ième pour le taux moyen & « \(\sqrt{1{,}44} = 1{,}2\) » \\
		\hline
		\texttt{COL} & Bon travail collectif, entraide efficace & Tous les membres participent \\
		\hline
		\texttt{VIT} & Rapidité d'exécution remarquable & Réponse donnée en moins de 30 s \\
		\hline
	\end{tabular}
\end{center}

\vspace{0.4cm}

\subsection*{Code abrégé -- Difficultés}

\begin{center}
	\renewcommand{\arraystretch}{1.4}
	\begin{tabular}{|c|>{\raggedright\arraybackslash}p{6cm}|>{\raggedright\arraybackslash}p{7.5cm}|}
		\hline
		\textbf{Symbole} & \textbf{Signification} & \textbf{Exemple observé} \\
		\hline
		\texttt{ADD} & Additionne encore les taux au lieu de multiplier & « \(-6+10-4=0\) » \\
		\hline
		\texttt{MOY} & Confond taux global et taux moyen & « moyenne des taux = taux moyen » \\
		\hline
		\texttt{REF} & Oublie la valeur de référence (\% sur quoi ?) & « \(+30\%\) puis \(-30\%\) s'annule » \\
		\hline
		\texttt{RET} & Oublie le coefficient de retour & « baisse de \(50\%\) \(\to\) baisse de \(50\%\) » \\
		\hline
		\texttt{RAC} & Bloque sur la racine n-ième ou la puissance \(1/n\) & « \((1+t)^5 = 1{,}27628\) non résolu » \\
		\hline
		\texttt{SGN} & Oublie d'interpréter le signe du résultat & « \(0{,}935\) » sans conclure \\
		\hline
		\texttt{CAL} & Erreur de calcul final alors que la démarche est juste & Bon CM, mauvais produit \\
		\hline
		\texttt{RED} & Rédaction absente ou trop allusive & Aucune phrase de conclusion \\
		\hline
		\texttt{BLQ} & Blocage total, abandon & Enveloppe rendue sans essai \\
		\hline
		\texttt{IND} & Un élève monopolise, les autres décrochent & Un seul écrit, les autres regardent \\
		\hline
	\end{tabular}
\end{center}

\vspace{0.4cm}

\subsection*{Grille d'observation express (à remplir pendant le jeu)}

\begin{center}
	{\small\itshape Notez 1 à 3 symboles par équipe et par manche. Utilisez le verso pour les cas particuliers.}
\end{center}

\vspace{0.2cm}

\begin{center}
	\renewcommand{\arraystretch}{1.7}
	\begin{tabular}{|>{\centering\arraybackslash}p{1.5cm}|>{\raggedright\arraybackslash}p{3.8cm}|>{\raggedright\arraybackslash}p{3.8cm}|>{\raggedright\arraybackslash}p{3.8cm}|>{\raggedright\arraybackslash}p{3.8cm}|}
		\hline
		\textbf{Éq.} & \textbf{Manche 1} & \textbf{Manche 2} & \textbf{Manche 3} & \textbf{Manches 4-5} \\
		\hline
		1 & & & & \\
		\hline
		2 & & & & \\
		\hline
		3 & & & & \\
		\hline
		4 & & & & \\
		\hline
	\end{tabular}
\end{center}

\vspace{0.4cm}

\subsection*{Fiche « Observations par équipe » (à remplir après la séance)}

\begin{center}
	{\small\itshape Section réservée à l'enseignant -- à compléter à partir des symboles notés pendant le jeu.}
\end{center}

\vspace{0.2cm}

\begin{center}
	\renewcommand{\arraystretch}{1.8}
	\begin{tabular}{|>{\centering\arraybackslash}p{1.2cm}|>{\raggedright\arraybackslash}p{4.4cm}|>{\raggedright\arraybackslash}p{4.4cm}|>{\raggedright\arraybackslash}p{4.4cm}|}
		\hline
		\textbf{Éq.} & \textbf{Points forts} & \textbf{Difficultés} & \textbf{À revoir} \\
		\hline
		1 & & & \\
		\hline
		2 & & & \\
		\hline
		3 & & & \\
		\hline
		4 & & & \\
		\hline
	\end{tabular}
\end{center}

\vspace{0.4cm}

\begin{conseil}
	\textbf{Comment remplir la fiche en 5 minutes ?}\\[0.2cm]
	\textbullet{} Regroupez les symboles identiques : plusieurs \texttt{ADD} pour l'équipe 2 \(\Rightarrow\) « Additionne encore les taux ».\\[0.2cm]
	\textbullet{} Traduisez chaque symbole en une phrase courte dans la bonne colonne (Point fort / Difficulté / À revoir).\\[0.2cm]
	\textbullet{} Une même notion peut apparaître en difficulté ET en « à revoir » : c'est normal, cela signifie qu'il faut la retravailler en priorité.\\[0.2cm]
	\textbullet{} Ne remplissez pas tout : 3 à 4 lignes par équipe suffisent pour un bilan de fin de séance.\\[0.2cm]
	\textbullet{} Servez-vous de la colonne « À revoir » pour préparer la séance suivante (exercices ciblés, rappels).
\end{conseil}

\vspace{0.4cm}

\subsection*{Exemple de fiche complétée}

\begin{center}
	\renewcommand{\arraystretch}{1.6}
	\begin{tabular}{|>{\centering\arraybackslash}p{1.2cm}|>{\raggedright\arraybackslash}p{4.4cm}|>{\raggedright\arraybackslash}p{4.4cm}|>{\raggedright\arraybackslash}p{4.4cm}|}
		\hline
		\textbf{Éq.} & \textbf{Points forts} & \textbf{Difficultés} & \textbf{À revoir} \\
		\hline
		1 & \texttt{CM}, \texttt{MULT}, \texttt{COL} : passage systématique par les coefficients. & \texttt{RAC} : blocage sur \((1+t)^5 = 1{,}27628\). & Racine n-ième ; taux moyen. \\
		\hline
		2 & \texttt{VF}, \texttt{INT} : bonne application de \(V_f = CM \times V_i\). & \texttt{ADD} : additionne encore les taux ; \texttt{REF} : oublie la valeur de référence. & Non-additivité des taux ; valeur de référence. \\
		\hline
		3 & \texttt{VIT}, \texttt{RAC} : rapide et à l'aise avec les racines. & \texttt{SGN} : oublie d'interpréter le signe. & Rédaction de la conclusion. \\
		\hline
		4 & \texttt{COL}, \texttt{CM} : bonne entraide, bases solides. & \texttt{RET} : oublie le coefficient de retour ; \texttt{BLQ} en Manche 5. & Coefficient de retour ; gestion du stress. \\
		\hline
	\end{tabular}
\end{center}
\end{document}